% ============================================================
% 75_pdf_identity.tex — PDF-Identitaet und Ownership-Nachweis
% ============================================================
% Alle dynamischen Felder (Titel, Release-ID, Build-Stamp, Git-Commit)
% werden vom Makefile via -pdflualatex="...\def\Foo{...}..." injiziert.
% Lokale Defaults via \providecommand sorgen dafür, dass das Dokument
% auch ohne Makefile (z.B. direkter latexmk-Aufruf) baubar bleibt.

% --- Owner (zentral hier setzen, einmal pro Person) ---
\newcommand{\ZSFOwnerDisplayName}{Loris Einar Kristj\'an Hliddal}
\newcommand{\ZSFOwnerNameASCII}{Loris Einar Kristjan Hliddal}
\newcommand{\ZSFOwnerID}{LEKH-ZSF-A1B2}

% --- Per Build injizierbare Felder (Defaults wenn ungesetzt) ---
\providecommand{\ZSFSubjectTitle}{ZSF Template}
\providecommand{\ZSFReleaseID}{DEV}
\providecommand{\ZSFReleaseLabel}{\ZSFReleaseID}
\providecommand{\ZSFBuildStamp}{manual}
\providecommand{\ZSFGitCommit}{nogit}
\newcommand{\ZSFBuildID}{\ZSFBuildStamp-\ZSFGitCommit}

% Sichtbarer dezenter Marker auf Seite 1 (klein, hellgrau)
\newif\ifZSFVisibleMarker
\ZSFVisibleMarkertrue

% XMP optional (graceful, wenn nicht installiert)
\IfFileExists{hyperxmp.sty}{\RequirePackage{hyperxmp}}{}
\RequirePackage{eso-pic}

\author{\ZSFOwnerNameASCII}
\title{\ZSFSubjectTitle}
\ifdefined\Author
  \Author{\ZSFOwnerNameASCII}
\fi
\ifdefined\Title
  \Title{\ZSFSubjectTitle}
\fi
\ifdefined\Subject
  \Subject{Originalfassung; owner=\ZSFOwnerID; release=\ZSFReleaseID}
\fi
\ifdefined\Keywords
  \Keywords{owner-name:\ZSFOwnerNameASCII, owner-id:\ZSFOwnerID, release-id:\ZSFReleaseID, build-id:\ZSFBuildID}
\fi

\hypersetup{
  pdftitle={\ZSFSubjectTitle},
  pdfauthor={\ZSFOwnerNameASCII},
  pdfsubject={Originalfassung; owner=\ZSFOwnerID; release=\ZSFReleaseID},
  pdfkeywords={owner-name:\ZSFOwnerNameASCII, owner-id:\ZSFOwnerID, release-id:\ZSFReleaseID, build-id:\ZSFBuildID},
  pdfcreator={LaTeX},
  pdfproducer={latexmk}
}

% Engine-Weiche: LuaTeX (>= 0.85) kennt \pdfinfo nicht mehr, sondern
% \pdfextension info; pdfTeX kennt nur \pdfinfo. Eine Quelle, beide Engines.
\ifdefined\pdfextension
  \protected\def\ZSFpdfinfo{\pdfextension info }
\else
  \let\ZSFpdfinfo\pdfinfo
\fi

\AtBeginDocument{
  % Low-level pdfinfo fallback (stabil bei Toolchains, die Author/Keywords sonst verwerfen).
  \ZSFpdfinfo{
    /Author (\ZSFOwnerNameASCII)
    /Title (\ZSFSubjectTitle)
    /Subject (Originalfassung; owner=\ZSFOwnerID; release=\ZSFReleaseID)
    /Keywords (owner-name:\ZSFOwnerNameASCII, owner-id:\ZSFOwnerID, release-id:\ZSFReleaseID, build-id:\ZSFBuildID)
  }
  \ifZSFVisibleMarker
    \AddToHook{shipout/firstpage}{
      \AtPageLowerLeft{
        \raisebox{\ZSFidentityMarkerInset}[0pt][0pt]{%
          \makebox[0pt][l]{%
            \hspace{\ZSFidentityMarkerInset}%
            {\color{IdentityMarkerColor}\ZSFfontIdentityMarker
            id:\ZSFOwnerID\space|\space rel:\ZSFReleaseID}%
          }%
        }%
      }%
    }
  \fi
}
